\documentclass[11pt,a4paper]{article}
\usepackage{amsmath,amssymb,amsthm}
\usepackage{listings}
\usepackage{xcolor}
\usepackage{geometry}
\geometry{margin=1in}

\lstset{
  language=Lean,
  basicstyle=\ttfamily\small,
  keywordstyle=\color{blue},
  commentstyle=\color{green!50!black},
  stringstyle=\color{red},
  breaklines=true,
  frame=single,
  numbers=left
}

\title{\textbf{Discrete Generative Contact Mechanics (dm³\_disc)}\\
A Lean-First Language Specification}
\author{Pablo Nogueira Grossi \\ G6 LLC / AXLE Project}
\date{v1.0 — April 2026}

\begin{document}

\maketitle

\begin{abstract}
This document presents dm³\_disc as a formally defined higher-order category that unifies the continuous dm³ framework (GCM Manifesto §2.2) with the Topographical Orthogonal Generative Theory (TOGT) operator algebra. The Collatz map is exhibited as a concrete object. Two axioms are proved; the remaining three are localized as precise analytic targets.
\end{abstract}

\section{The Discrete dm³ Grammar (Lean-first)}

\subsection{Syntax}
\begin{lstlisting}
structure TOGTGrammar where
  C : ℕ → ℕ
  K : ℕ → Bool
  F : ℕ → ℕ
  U : ℕ → ℕ
\end{lstlisting}
C = compression, K = constraint (curvature), F = folding, U = unfolding.

\subsection{Composition}
The generative operator is the composite \(G = U \circ F \circ K \circ C\).

\section{The Discrete dm³ Category}

\subsection{Objects}
\begin{lstlisting}
structure DiscreteDm3System where
  X        : Type
  T        : X → X
  V        : X → ℝ
  Phi      : X → ℝ
  Gamma    : Finset X
  grammar  : TOGTGrammar
  -- eight axioms as fields (see v1.6)
\end{lstlisting}
An object is a countable discrete state space equipped with a generative map \(T\), contact structure, Lyapunov function, and the TOGT grammar.

\subsection{Morphisms}
\begin{lstlisting}
structure DiscreteDm3Hom (A B : DiscreteDm3System) where
  f        : A.X → B.X
  map_T    : ∀ x, f (A.T x) = B.T (f x)
  map_Phi  : ∀ x, |B.Phi (f x) - A.Phi x| ≤ 1
  map_V    : ∀ x, B.V (f x) ≤ A.V x + 1
\end{lstlisting}
Morphisms are
